\documentclass[11pt]{article}

\usepackage[margin=1in]{geometry}
\usepackage{amsmath,amssymb,amsthm}
\usepackage[T1]{fontenc}
\usepackage[hidelinks]{hyperref}

\newtheorem{theorem}{Theorem}
\newtheorem{lemma}{Lemma}
\newtheorem{corollary}{Corollary}

\newcommand{\diam}{\operatorname{diam}}
\newcommand{\limsupn}{\limsup_{n\to\infty}}

\title{A Bounded-Domain Collapse for Asymptotically Holder Nonexpansive Type Maps}
\author{Run \texttt{fa\_banach\_001}, agent \texttt{agent\_super\_00}}
\date{2026-07-01}

\begin{document}
\maketitle

\section*{Source And Target}

The source paper is C. S. Barroso and C. S. R. da Silva,
\emph{Fixed point results for asymptotically Holder nonexpansive type
mappings}, arXiv:2605.30678.  After proving a Goebel--Kirk style theorem for
asymptotically Holder nonexpansive type mappings under the hypothesis
\(\epsilon_0(X)<1\), the paper says that the extent to which other classical
results, including those of Dominguez Benavides and Lorenzo Ramirez, extend to
this framework remains a subject for future investigation.

The supporting source is T. Dominguez Benavides and P. Lorenzo Ramirez,
\emph{Fixed points for mappings of asymptotically nonexpansive type},
arXiv:2201.02802, Fixed Point Theory 24 (2023), 569--582.  In particular,
their Theorem 5.3 states that if \(X\) has uniform normal structure, \(C\) is a
bounded convex subset of \(X\), and \(T:C\to C\) is asymptotically
nonexpansive type, then there exists \(x\in C\) such that \(T^n x\to x\).

\section*{Definitions}

Let \(C\) be a bounded subset of a Banach space \(X\). A map \(T:C\to C\) is
of asymptotically Holder nonexpansive type if there is a sequence
\(\alpha_n\in(0,1]\) with \(\alpha_n\to 1\) such that, for every \(x\in C\),
\[
\limsupn \sup_{y\in C}\left(\|T^n x-T^n y\|-\|x-y\|^{\alpha_n}\right)\le 0.
\]
The ordinary asymptotically nonexpansive type condition is the same assertion
with \(\|x-y\|\) in place of \(\|x-y\|^{\alpha_n}\).

\section*{Proof Intuition}

The Holder version looks weaker at small scales, because \(t^\alpha\ge t\) on
\([0,1]\) when \(0<\alpha\le 1\).  But on any bounded interval the functions
\(t^\alpha\) converge uniformly to \(t\) as \(\alpha\to 1\).  Therefore the
Holder-type inequality differs from the usual asymptotically nonexpansive type
inequality by an additive error that tends to zero uniformly over the whole
domain.  Once this reduction is made, the fixed-point machinery already known
for asymptotically nonexpansive type maps applies verbatim.

\begin{lemma}[Bounded-domain collapse]\label{lem:collapse}
Let \(C\) be a bounded subset of a Banach space \(X\), and let
\(\alpha_n\in(0,1]\) satisfy \(\alpha_n\to 1\). If \(T:C\to C\) is of
asymptotically Holder nonexpansive type with exponent sequence
\((\alpha_n)\), then \(T\) is of ordinary asymptotically nonexpansive type.
\end{lemma}

\begin{proof}
Put \(D=\diam(C)<\infty\).  For \(n\ge 1\), define
\[
 \eta_n=\sup_{0\le t\le D}\bigl(t^{\alpha_n}-t\bigr).
\]
The sequence \((\alpha_n)\) is eventually contained in some compact interval
\([\alpha_0,1]\subset(0,1]\).  The map \((t,\alpha)\mapsto t^\alpha\) is
continuous on \([0,D]\times[\alpha_0,1]\), hence uniformly continuous there;
therefore \(\eta_n\to 0\).

Fix \(x\in C\).  For every \(y\in C\) and every \(n\),
\[
\|T^n x-T^n y\|-\|x-y\|
\le
\left(\|T^n x-T^n y\|-\|x-y\|^{\alpha_n}\right)+\eta_n .
\]
Taking the supremum over \(y\in C\) and then the limsup over \(n\) gives
\[
\limsupn \sup_{y\in C}\left(\|T^n x-T^n y\|-\|x-y\|\right)
\le 0+0=0.
\]
This is the ordinary asymptotically nonexpansive type condition.
\end{proof}

\begin{theorem}[Uniform-normal-structure transfer]\label{thm:uns}
Let \(X\) be a Banach space with uniform normal structure, let \(C\subset X\)
be a bounded convex subset, and let \(T:C\to C\) be of asymptotically Holder
nonexpansive type. Then there exists \(x\in C\) such that \(T^n x\to x\).
Consequently, under any continuity hypothesis that turns such a convergent
orbit into a fixed point, for instance if some fixed iterate \(T^N\) is
continuous, \(T\) has a fixed point.
\end{theorem}

\begin{proof}
By Lemma~\ref{lem:collapse}, \(T\) is asymptotically nonexpansive type in the
ordinary sense.  Dominguez Benavides--Lorenzo Ramirez Theorem 5.3 therefore
applies and yields \(x\in C\) with \(T^n x\to x\).

If \(T^N\) is continuous, then \(T^{n+N}x=T^N(T^n x)\to T^Nx\), while also
\(T^{n+N}x\to x\).  Hence \(T^N x=x\).  The orbit of \(x\) is then periodic
with period dividing \(N\), but the full orbit \(T^n x\) converges to \(x\);
therefore all points in the finite cycle equal \(x\), in particular \(Tx=x\).
\end{proof}

\begin{corollary}
Every bounded-domain fixed-point or orbit-convergence theorem whose hypothesis
is the ordinary asymptotically nonexpansive type condition immediately applies
to asymptotically Holder nonexpansive type mappings with \(\alpha_n\to 1\).
In particular, the 2026 theorem under \(\epsilon_0(X)<1\) is part of a larger
uniform-normal-structure transfer.
\end{corollary}

\section*{Limitations And Novelty Check}

This packet does not settle the Dowling--Lennard--Turett renorming question
for spaces containing \(\ell_1\); the source paper explicitly separates that
general asymptotically nonexpansive problem from its Holder-class result.
The result here addresses the bounded-domain Holder-type framework and the
source paper's broader invitation to extend classical asymptotic-type fixed
point theorems.

The cheap run indexes were searched for arXiv:2605.30678 and for combinations
of ``Holder'', ``asymptotically nonexpansive type'', and ``uniform normal
structure''.  No duplicate packet was found.  A local full-source search found
the supporting 2022/2023 theorem but no packet recording this bounded-domain
collapse for the 2026 Holder-type paper.

\begin{thebibliography}{9}
\bibitem{BarrosoSilva2026}
C. S. Barroso and C. S. R. da Silva,
\emph{Fixed point results for asymptotically Holder nonexpansive type
mappings}, arXiv:2605.30678, 2026.

\bibitem{DBLR2023}
T. Dominguez Benavides and P. Lorenzo Ramirez,
\emph{Fixed points for mappings of asymptotically nonexpansive type},
Fixed Point Theory 24 (2023), 569--582; arXiv:2201.02802.
\end{thebibliography}

\end{document}
